\NativeTitle{49}{第六单元素养评价（A）}{（小数的意义与性质）}{90}{5}{一 & 二 & 三 & 四 & 五}{ &  &  &  & }
\NativeSlot{275}{59}{%
\ExamGroup{一、填空。{\normalfont\kaishu\fontsize{10.8bp}{14bp}\selectfont （每空1分，共24分）}}

% p049_s01

\ExamQuestion{1.}{5}%
中国国家博物馆遵循“开门办馆”的理念，分期分批对社会公布馆藏藏品目录，第一期共公布藏品约30.64万件。横线上的小数读作\mbox{（\hspace{7em}）}，这个小数的计数单位是\mbox{（\hspace{4.5em}）}，它有\mbox{（\hspace{1.8em}）}个这样的计数单位。\par
}
\NativeSlot{374.0}{41}{%
% p049_s02

\ExamQuestion{2.}{5}%
一个小数由8个0.1和7个0.01组成，这个小数是\mbox{（\hspace{1.8em}）}，再加上\mbox{（\hspace{1.8em}）}个\(\frac1{100}\)就是1。\par
}
\NativeSlot{444.0}{42}{%
% p049_s03

\ExamQuestion{3.}{5}%
把168000改写成用“万”作单位的数是\mbox{（\hspace{1.8em}）}万，保留整数约是\mbox{（\hspace{1.8em}）}万；把995000000改写成以“亿”为单位的数是\mbox{（\hspace{1.8em}）}亿，精确到十分位约是\mbox{（\hspace{1.8em}）}亿。\par
}
\NativeSlot{518}{36}{%
% p049_s04

\ExamQuestion{4.}{5}%
3分米＝\mbox{（\hspace{1.8em}）}米；5.06吨＝\mbox{（\hspace{1.8em}）}吨\mbox{（\hspace{1.8em}）}千克；32克＝\mbox{（\hspace{1.8em}）}千克；5平方米3平方分米＝\mbox{（\hspace{1.8em}）}平方米。\par
}
\NativeSlot{568}{58}{%
% p049_s05

\ExamQuestion{5.}{5}%
把0.504、0.45、0.045、0.54、0.405、0.054按从小到大的顺序排列是\mbox{（\hspace{1.8em}）}。\par

% p049_s06

\ExamQuestion{6.}{5}%
元元在读一个小数时，把小数点漏读了，结果读成八万四千零二。若原来的小数要读一个零，则原来的小数是\mbox{（\hspace{1.8em}）}。\par
}
\NativeSlot{671}{65}{%
% p049_s07

\ExamQuestion{7.}{5}%
在圆圈里填上“＞”“＜”或“＝”。\(6.02\ \bigcirc\ 6.002\)，\allowbreak{}\(3.5\ \bigcirc\ 3.50\)，9千克50克\(\bigcirc\)9.5千克；\allowbreak{}\(7.91\ \bigcirc\ 7.97\)，0.57 m\(\bigcirc\)57 cm，\allowbreak{}\(8.26\ \bigcirc\ 4.62\)。\par
}
